% ══════════════════════════════════════════════════════════════════
%  FICHE D'EXERCICES — Chapitre 8 : Les amines
% ══════════════════════════════════════════════════════════════════
\section{Exercices type Baccalauréat}
\textit{Données : $M(\ce{H})=1$, $M(\ce{C})=12$, $M(\ce{N})=14$ \si{\gram\per\mol} ;
$\ce{H3O+}$, $pK_e=14$ à \SI{25}{\celsius}.}

\rubrique{Classes, nomenclature}

\exo{}\diff{1} Donner le nom et la classe (primaire, secondaire, tertiaire) des
amines : a)~$\ce{CH3-NH2}$ ; b)~$\ce{CH3-NH-CH3}$ ; c)~$\ce{(CH3)3N}$ ;
d)~$\ce{CH3-CH2-NH2}$.

\exo{}\diff{1} Écrire les formules semi-développées de toutes les amines de formule
brute $\ce{C3H9N}$, les nommer et préciser leur classe.

\rubrique{Caractère basique}

\exo{}\diff{1} Écrire l'équation de la réaction de la méthylamine $\ce{CH3-NH2}$
avec l'eau. Identifier le couple acide/base mis en jeu.

\exo{}\diff{1} Écrire l'équation de la réaction de l'éthylamine avec l'acide
chlorhydrique. Nommer le sel formé.

\exo{}\diff{2} Une solution d'éthylamine de concentration
$C = \SI{0.10}{\mol\per\litre}$ a un pH de $11{,}8$.
\begin{enumerate}[label=\arabic*)]
\item Écrire l'équation de la réaction de l'éthylamine avec l'eau.
\item Montrer que l'éthylamine est une base faible.
\item Calculer les concentrations des espèces et en déduire le $pK_a$ du couple
$\ce{C2H5-NH3+}/\ce{C2H5-NH2}$.
\end{enumerate}

\exo{}\diff{2} Comparer la basicité de l'ammoniac et de la méthylamine. Quelle est
l'influence du groupe alkyle ? (On pourra s'appuyer sur les $pK_a$ : $\ce{NH4+}/\ce{NH3}$,
$pK_a = 9{,}2$ ; $\ce{CH3NH3+}/\ce{CH3NH2}$, $pK_a = 10{,}7$.)

\rubrique{Détermination de formule}

\exo{}\diff{2} Une amine primaire saturée $A$ a une masse molaire
$M = \SI{59}{\gram\per\mol}$.
\begin{enumerate}[label=\arabic*)]
\item Déterminer sa formule brute (formule générale $\ce{C_nH_{2n+3}N}$).
\item Écrire les isomères amines primaires et les nommer.
\end{enumerate}

\exo{}\diff{3} La combustion complète de \SI{4.5}{\gram} d'une amine saturée $A$
fournit \SI{8.8}{\gram} de $\ce{CO2}$, \SI{6.3}{\gram} d'eau et du diazote. La masse
molaire de $A$ est \SI{45}{\gram\per\mol}.
\begin{enumerate}[label=\arabic*)]
\item Déterminer la formule brute de $A$.
\item Donner les isomères de $A$, leurs noms et leurs classes.
\item $A$ réagit avec l'acide chlorhydrique. Écrire l'équation pour l'isomère amine
primaire à chaîne droite.
\end{enumerate}

\rubrique{Problème de synthèse — type Baccalauréat}

\sujetbac[d'après Baccalauréat C, Cameroun]
On considère une amine $A$ de formule brute $\ce{C2H7N}$.
\begin{enumerate}[label=\arabic*)]
\item Écrire les formules semi-développées des isomères de $A$, les nommer et donner
leur classe.
\item On dissout l'un des isomères (amine primaire) dans l'eau ; la solution obtenue
est basique. Écrire l'équation de la réaction avec l'eau et nommer les espèces.
\item On dose \SI{20}{\milli\litre} de cette solution d'amine
($C = \SI{0.10}{\mol\per\litre}$) par une solution d'acide chlorhydrique
$C_a = \SI{0.10}{\mol\per\litre}$. Écrire l'équation du dosage et calculer le volume
d'acide versé à l'équivalence.
\end{enumerate}
